\section{Cross-dataset transfer benchmark}
\label{imp:app:openml}
\label{imp:sec:task-openml}
\label{imp:sec:results-openml}

\subsection{Design}


To test behavior beyond economic microdata, we run the method surface over
the six OpenML AutoML Benchmark regression datasets used in the
\microimpute{} manuscript's appendix \citep{vanschoren2014openml} ---
\texttt{space\_ga}, \texttt{elevators}, \texttt{brazilian\_houses},
\texttt{onlinenewspopularity}, \texttt{abalone}, and \texttt{house\_sales}
--- treating each dataset's target as the imputed column under uniform
weights.


\subsection{Results}


\begin{table}[t]
\centering
\footnotesize
\resizebox{\textwidth}{!}{% Generated by `imp tables` -- do not edit by hand. Wasserstein-1 to donor, mean over 10 paired seeds; lower is better.
\begin{tabular}{lrrrrrrr}
\hline
Method & abalone & brazilian\_houses & elevators & house\_sales & onlinenewspopularity & space\_ga & Mean rank \\
\hline
populace-fit & 0.254 & 166.64 & 0.000 & 7,080 & 274.67 & 0.008 & 3.83 \\
\quad -- unweighted & 0.257 & 184.92 & 0.000 & 6,744 & 233.10 & 0.008 & \textbf{3.33} \\
\quad -- unchained & 0.254 & 166.64 & 0.000 & 7,080 & 274.67 & 0.008 & 3.83 \\
\quad -- ungated/unchained forest & 0.174 & 191.36 & 0.000 & 8,461 & 546.16 & 0.012 & 5.00 \\
OLS & 0.532 & 212.09 & 0.002 & 90,518 & 6,117 & 0.015 & 8.83 \\
Quantile regression & 0.321 & 212.04 & 0.001 & 24,115 & 382.98 & 0.085 & 7.83 \\
NND hot deck (py-statmatch) & 0.176 & 199.12 & 0.000 & 10,212 & 161.43 & 0.010 & 4.00 \\
Weighted marginal draw & 0.148 & 210.50 & 0.000 & 7,135 & 143.52 & 0.011 & \textbf{3.33} \\
\hline
\end{tabular}
}
\caption{Wasserstein-1 to the donor distribution across the six OpenML
regression datasets (uniform weights), mean over ten paired seeds, with mean
rank across datasets.}
\label{imp:tab:openml}
\end{table}

Table~\ref{imp:tab:openml} runs the surface over the six OpenML regression
datasets under uniform weights. Two observations. First, with weights
uninformative by construction, the candidate and its unweighted ablation are
statistically the same method, and they rank near the top (mean ranks 3.8 and
3.3) alongside the unconditional marginal draw (3.3); the linear methods rank
last (7.8 and 8.8). Second --- and more important for this paper's argument --- the best
mean rank on a \emph{marginal} distance is achieved by a method with no
conditional structure at all. Marginal metrics saturate: they reward drawing
from the right marginal and cannot rank what imputation is actually for.
The population-view harness exists because of exactly this ceiling.

